\chapter{Introduction}
\label{chap:introduction}

\section{Pr\'esentation de l'application}

\textbf{CGP -- Annuaire Formation} est une application web de gestion de l'annuaire universitaire. Elle centralise les informations sur les structures acad\'emiques, les responsables p\'edagogiques, les formations et le personnel de secr\'etariat de l'\'etablissement.

Ses principales fonctionnalit\'es sont\,:

\begin{itemize}
  \item \textbf{Annuaire et recherche} : consultation multi-crit\`eres des responsables, formations, structures et secr\'etariats.
  \item \textbf{Gestion des affectations} : association de personnes \`a des r\^oles au sein des structures.
  \item \textbf{Organigrammes} : g\'en\'eration et consultation d'organigrammes hi\'erarchiques (vue structures ou vue personnes).
  \item \textbf{Gestion multi-ann\'ees} : cr\'eation, clonage, archivage et suppression d'ann\'ees acad\'emiques.
  \item \textbf{Import et export} : \'echanges de donn\'ees via des classeurs Excel standardis\'es.
  \item \textbf{D\'el\'egations} : transfert temporaire de droits entre utilisateurs.
  \item \textbf{Signalements} : d\'eclaration et suivi d'erreurs sur les donn\'ees.
  \item \textbf{Journal d'audit} : tra\c{c}abilit\'e compl\`ete des actions effectu\'ees dans l'application.
\end{itemize}

\screenshot{img_accueil}{Page d'accueil de l'application CGP}{Capturer la page d'accueil ou de connexion de l'application telle qu'elle appara\^it dans le navigateur web (Google Chrome ou Firefox recommand\'e). L'URL de l'application doit \^etre visible dans la barre d'adresse. Utiliser une fen\^etre en plein \'ecran, r\'esolution 1920$\times$1080.}

\section{Architecture g\'en\'erale}

L'application est organis\'ee autour de concepts fondamentaux\,:

\begin{description}[style=nextline, leftmargin=3cm, font=\bfseries]
  \item[Ann\'ee acad\'emique] Unit\'e de temps principale. Toutes les donn\'ees (structures, affectations, organigrammes) sont rattach\'ees \`a une ann\'ee. Une seule ann\'ee peut \^etre \og{} EN\_COURS \fg{} \`a la fois.
  \item[Structure (entit\'e)] Unit\'e organisationnelle : Composante, D\'epartement, Mention, Parcours, Niveau, Service, etc.
  \item[Utilisateur] Toute personne ayant un acc\`es \`a l'application.
  \item[Affectation] Lien entre un utilisateur, un r\^ole et une structure pour une ann\'ee donn\'ee.
  \item[R\^ole] Fonction occupée par un utilisateur dans une structure (ex. : Directeur de composante, Responsable de formation).
  \item[Organigramme] Repr\'esentation graphique hi\'erarchique g\'en\'er\'ee \`a partir des affectations.
  \item[D\'el\'egation] Transfert temporaire de droits d'un utilisateur vers un autre.
  \item[Signalement] Rapport d'erreur soumis par n'importe quel utilisateur.
\end{description}

\section{Acc\`es \`a l'application}

L'application CGP est accessible via un navigateur web \`a l'adresse communiqu\'ee par votre administrateur syst\`eme.

\begin{notebox}[Navigateurs compatibles]
L'application est optimis\'ee pour\,: \textbf{Google Chrome} (recommand\'e), \textbf{Mozilla Firefox} et \textbf{Microsoft Edge}. \\
Internet Explorer et les navigateurs anciens ne sont \textbf{pas} support\'es.
\end{notebox}

\begin{notebox}[Pr\'erequis r\'eseau]
  Une connexion au r\'eseau de l'universit\'e est n\'ecessaire (connexion filaire, Wi-Fi universit\'e ou VPN selon la configuration de votre \'etablissement).
\end{notebox}

\section{Organisation du manuel}

Ce manuel est structur\'e comme suit\,:

\medskip
\begin{tabularx}{\textwidth}{C{2.5cm}L{4cm}X}
  \toprule
  \rowcolor{cgplightblue}
  \textbf{Chapitre} & \textbf{Titre} & \textbf{Contenu} \\
  \midrule
  Chapitre~1 & Introduction & Ce chapitre -- pr\'esentation g\'en\'erale \\
  \rowcolor{cgplightgray}
  Chapitre~2 & Connexion et profil & Authentification, profil utilisateur, notifications \\
  Chapitre~3 & R\^oles et droits & Pr\'esentation de tous les r\^oles, matrice des droits \\
  \rowcolor{cgplightgray}
  Chapitre~4 & Services Centraux & Manuel complet pour le r\^ole \texttt{services-centraux} \\
  Chapitre~5 & Directeur de Composante & Manuel pour le r\^ole \texttt{directeur-composante} \\
  \rowcolor{cgplightgray}
  Chapitre~6 & Directions et Responsables & Manuel pour les r\^oles de direction (d\'epartement, mention, etc.) \\
  Chapitre~7 & Secr\'etariat P\'edagogique & Manuel pour \texttt{secretariat-pedagogique} \\
  \rowcolor{cgplightgray}
  Chapitre~8 & Utilisateur Standard & Manuel pour \texttt{utilisateur-simple} et \texttt{lecture-seule} \\
  Chapitre~9 & Fonctionnalit\'es communes & Annuaire, profil, signalements, organigrammes (lecture) \\
  \bottomrule
\end{tabularx}

\medskip
\begin{tipbox}[Comment utiliser ce manuel ?]
  Si vous cherchez la proc\'edure pour une fonctionnalit\'e sp\'ecifique, consultez d'abord le chapitre correspondant \`a votre r\^ole. Si la fonctionnalit\'e n'y est pas d\'etaill\'ee, reportez-vous au chapitre~\ref{chap:commun} (Fonctionnalit\'es communes).
\end{tipbox}

\section{Conventions typographiques}

Les conventions utilis\'ees dans ce manuel sont les suivantes\,:

\medskip
\begin{tabularx}{\textwidth}{C{3cm}X}
  \toprule
  \rowcolor{cgplightblue}
  \textbf{Convention} & \textbf{Signification} \\
  \midrule
  \textbf{Gras} & \'El\'ement d'interface : bouton, menu, onglet (ex. : \textbf{Cr\'eer}, \textbf{Valider}) \\
  \rowcolor{cgplightgray}
  \textit{Italique} & Nom de champ de formulaire (ex. : \textit{Nom}, \textit{Courriel institutionnel}) \\
  \texttt{monospace} & Valeur technique, code, nom de r\^ole (ex. : \texttt{services-centraux}) \\
  \rowcolor{cgplightgray}
  \faInfoCircle & Note informative compl\'ementaire \\
  \faExclamationTriangle & Avertissement -- point important \`a ne pas n\'egliger \\
  \rowcolor{cgplightgray}
  \faLightbulb & Conseil -- bonne pratique recommand\'ee \\
  \faUserShield & Indication du ou des r\^oles concern\'es par la section \\
  \bottomrule
\end{tabularx}
